% Concept note: Interactive Economy Simulator for Economics Education and Research
% Compile with:  tectonic concept-note.tex   (or xelatex / pdflatex)
\documentclass[10pt,a4paper]{article}

% ---------- Edit these ----------
\newcommand*{\authorname}{Dr Sulaxan}
\newcommand*{\contactline}{}          % e.g. name@example.com \textperiodcentered{} +91 98xxx xxxxx
\newcommand*{\institute}{}            % e.g. Department of Economics, XYZ University
% --------------------------------

\usepackage[a4paper,top=1.1cm,bottom=1.1cm,left=1.5cm,right=1.5cm]{geometry}
\usepackage{tgpagella}
\usepackage[T1]{fontenc}
\usepackage{microtype}
\usepackage{xcolor}
\usepackage{tabularx,booktabs}
\usepackage{enumitem}
\usepackage{multicol}
\usepackage{titlesec}
\usepackage[hidelinks]{hyperref}

\definecolor{ink}{HTML}{14213D}
\definecolor{accent}{HTML}{C8611A}

\pagestyle{empty}
\setlength{\parindent}{0pt}
\setlength{\parskip}{2pt}
\setlength{\multicolsep}{2pt}
\setlength{\columnsep}{18pt}
\titleformat{\section}{\fontsize{11.5}{13}\selectfont\bfseries\color{ink}}{}{0pt}{}
\titlespacing*{\section}{0pt}{5pt}{1pt}
\setlist{nosep,leftmargin=1.2em,topsep=1pt}
\renewcommand{\arraystretch}{1.05}

\begin{document}
\fontsize{9.6}{11.5}\selectfont

{\Large\bfseries\color{ink} Concept Note: Interactive Economy Simulator\\[2pt]
for Economics Education and Research\par}
\vspace{4pt}
{\color{accent}\rule{\linewidth}{1.2pt}}\par
\vspace{2pt}
\textbf{\authorname}%
\ifx\institute\empty\else\ \textperiodcentered\ \institute\fi
\ifx\contactline\empty\else\ \textperiodcentered\ \contactline\fi
\hfill 17 November 2026
\vspace{4pt}

We propose to build and pilot, with the Institute, an interactive simulator of the Indian economy for teaching and research. Students and faculty set policy and watch a living economy respond.

\section{The problem}
\begin{itemize}
  \item \textbf{Macroeconomics is taught on static diagrams.} IS-LM, the multiplier, the Phillips curve and Solow are taught as chalkboard models. Students rarely see how an economy actually responds to a policy over time.
  \item \textbf{Simulation is missing from current syllabi.} In the BSc and MSc Economics syllabi we reviewed, agent-based modelling is not mentioned once, and simulation appears only twice, never in macroeconomics. Yet every programme now has R, Python or Excel labs.
  \item \textbf{No digital tool fits the Indian context.} Existing simulators are built around the US Federal Reserve or the European Central Bank. In India, teaching relies on role-plays, token games and paper-trading apps that cover stocks, not the economy. None models the RBI's repo rate, CRR and SLR, GST, the informal sector or monsoon-driven food inflation.
\end{itemize}

\section{Proposed solution}
A browser-based \textbf{agent-based economy simulator}. Thousands of simulated households, firms and banks, plus a government and the RBI, interact month by month. Inflation, unemployment, inequality and crises emerge from their behaviour rather than from fixed formulas.
\begin{itemize}
  \item \textbf{Textbook mode} shows the IS-LM, AD-AS or multiplier result beside what the simulated economy actually does, and explains the gap. It strengthens exam topics instead of replacing them.
  \item \textbf{India scenarios:} \emph{You are the RBI Governor} (repo, CRR, SLR, OMO), \emph{Finance Minister's Budget} (GST, deficit), \emph{Bad monsoon}, \emph{Oil price spike and the rupee}.
  \item \textbf{One engine for every user:} a mobile-friendly web app for classrooms, and a Python package for researchers. Results can be reproduced exactly and shared as links.
\end{itemize}

\section{Who benefits and how it fits the syllabus}
{\footnotesize
\begin{tabularx}{\linewidth}{@{}>{\raggedright\arraybackslash}p{0.31\linewidth}X>{\raggedright\arraybackslash}p{0.11\linewidth}@{}}
\toprule
\textbf{Course} & \textbf{How the simulator is used} & \textbf{Users} \\
\midrule
Intro and Intermediate Macroeconomics & Multiplier, IS-LM, AD-AS and the Phillips curve shown live next to the textbook model & BSc \\
Money and Banking & Credit creation; repo, CRR and SLR in the RBI Governor scenario & BSc, MSc \\
Public Finance & Taxes, spending and the deficit; crowding out & BSc, MSc \\
Indian Economy & Calibrated shocks: monsoon, oil and the rupee, COVID & BSc, MSc \\
Advanced Macro and Monetary Economics & Taylor rule; comparing the simulation with DSGE models; expectations & MSc \\
R and Python labs, Econometrics & Scripted experiments; simulated data where the true effect is known & BSc, MSc \\
Dissertations and research & A ready platform for policy experiments that others can reproduce & MSc, faculty \\
\bottomrule
\end{tabularx}\par}

\section{Key features}
\begin{multicols}{2}
\begin{itemize}
  \item Scenarios labelled by syllabus unit, each with a teacher kit: worksheet, questions and answer guide.
  \item An ``explain why'' panel that traces each outcome back to the actual flows in the economy.
  \item Switchable schools of thought (Keynesian, monetarist and others), so students can see how assumptions change conclusions.
  \item Calibration to RBI and MOSPI data, and checks against real-world patterns such as the Phillips curve and Okun's law.
  \item Runs on low-cost Android phones; English and Hindi first.
  \item An open-source engine whose models are documented in the standard ODD format, so faculty can audit and extend it.
\end{itemize}
\end{multicols}

\noindent\begin{minipage}[t]{0.485\linewidth}
\section{What we request from the Institute}
\begin{itemize}
  \item \textbf{Faculty advisors:} one or two who teach Macroeconomics or Money and Banking, to review the scenarios and validate the model.
  \item \textbf{A pilot:} one semester in one or two courses, with a pre-test and post-test on concept understanding.
  \item \textbf{Data and research:} guidance on Indian calibration data, and a joint working paper or case study from the pilot.
  \item \textbf{No cost to the Institute during the pilot.} Free access for participating students and faculty, with acknowledgement as a founding academic partner.
\end{itemize}
\end{minipage}\hfill
\begin{minipage}[t]{0.485\linewidth}
\section{Timeline and expected outcomes}
\begin{enumerate}
  \item \textbf{Months 0--3, working demo:} the core economy, the RBI Governor and Budget scenarios, and textbook mode, reviewed with faculty.
  \item \textbf{Months 4--8, classroom pilot:} one semester in partner courses, with teacher kits and a measured learning outcome.
  \item \textbf{Months 9--12, research release:} the Python package, calibration to Indian data, and a joint paper or case study.
\end{enumerate}
\vspace{2pt}
\textbf{Outcomes:} better student understanding of macro and monetary topics, a hands-on bridge between theory and the R and Python labs, and a citable India-calibrated research tool. The Institute is recognised as a pioneer in simulation-based economics teaching.
\end{minipage}

\end{document}
